\documentclass[11pt]{article}
\usepackage[margin=1in]{geometry}
\usepackage{amsmath,amssymb,amsthm}
\usepackage{hyperref}
\hypersetup{hidelinks}
\usepackage{booktabs}

\newtheorem{theorem}{Theorem}
\newtheorem{lemma}[theorem]{Lemma}
\newtheorem{corollary}[theorem]{Corollary}
\newtheorem{proposition}{Proposition}
\theoremstyle{definition}
\newtheorem{definition}{Definition}
\theoremstyle{remark}
\newtheorem*{remark}{Remark}

\setlength{\parskip}{4pt plus 1pt minus 1pt}
\setlength{\parindent}{0pt}

\title{Bounded Context Architecture:\\Graph Structure, Escalation, Normalization,\\and What Persists}
\author{Patrick D.\ McCarthy}
\date{}

\begin{document}
\maketitle

\begin{abstract}
Any system that accumulates knowledge under bounded active context
faces a fundamental tradeoff between consistency and availability.
We show that this single constraint---bounded context---predicts
four architectural consequences that independently emerge across
software, biological, and organizational systems.

First, a directed graph with typed edges is the cost-optimal
representation: factoring shared structure preserves information
while reducing context load, and sub-linear dependency queries
require stored adjacency. Second, as action repertoire grows,
control architectures converge toward escalation: lower levels
resolve locally, signaling upward only on failure. Third, shared
propositions precipitate from the action space into a separated
knowledge structure---normalization in Codd's sense---eliminating
the update anomalies that manifest as catastrophic forgetting in
architectures without partition boundaries. Fourth, agency
($\gamma > 0$) and learning efficiency ($\eta_M > 0$) are the
minimal structural invariants of persistence under positive novelty
rate; their product $I(E) = \gamma \cdot \eta_M$ is a natural
measure of intelligence as rate of adaptation to novelty.

These results are not independent insights but consequences of one
constraint. Empirical confirmation comes from cancer genomics:
grouping mutations by coupling-channel boundaries (the graph's
partition structure) predicts overall survival with $p < 10^{-14}$
across 71,983 patients, while total mutation count does not.
\end{abstract}

% ==================================================================
% SECTION 1: INTRODUCTION
% ==================================================================
\section{Introduction}
\label{sec:intro}

% TODO: Merge the best introductory material from Papers 0-4.
% Key elements to include:
% - The bounded context constraint (from Paper 1)
% - The two strategies: flat/fungible vs structured (from Paper 1)
% - The crossover argument (from Paper 1)
% - Physarum polycephalum existence proof (from Paper 1)
% - "The industry is slime-molding" (from Paper 1)
% - The single-constraint unification claim (new, the merger's contribution)
% - The software-as-graph enumeration (from Paper 0, condensed)
% - Cancer as defection from graph back to swarm (one sentence)
% - Brief roadmap of the four results

% PLACEHOLDER: Introduction text to be assembled from source papers.
% See merger notes in publications/merge_notes.md

% ==================================================================
% SECTION 2: FORMAL FRAMEWORK
% ==================================================================
\section{Formal Framework}
\label{sec:framework}

% TODO: One unified set of definitions, appearing once.
% From Paper 1: World State Space, Proposition/Knowledge Base,
%   Bounded Active Context, Growth/Curation Regimes,
%   Contextual Retrieval, Factoring, Information Preservation
% From Paper 2: Encoding Hierarchy, Delegation Overhead, Escalation Policy
% From Paper 3: Action Space, Context Load, Entity
% From Paper 4: Agency, Action Feedback, Information Gain,
%   Novelty Rate, Free Context, Higher-Order Function Space M,
%   M Efficiency, Intelligence, Persistence
%
% Selective Necessity remark appears once here.

% PLACEHOLDER: Definitions to be assembled from source papers.

% ==================================================================
% SECTION 3: GRAPH STRUCTURE
% ==================================================================
\section{Graph Structure Is Cost-Optimal Under Bounded Context}
\label{sec:graph}

% TODO: From Paper 1:
% - Proposition 1 (Graph Optimality from Information Preservation)
%   with the five-option exhaustive argument
% - Theorem 1 (Graph Structure from Retention Dynamics)
%   with the O(n) lower bound proof
% - Corollary: Edge Growth Dominates Node Growth
% - Theorem: Contextual Retrieval Requires Adjacency Computation
% - The attention/transformer convergence argument
% - Catastrophic forgetting as empirical confirmation

% PLACEHOLDER: Graph section to be assembled from Paper 1.

% ==================================================================
% SECTION 4: ESCALATION
% ==================================================================
\section{Convergence to Escalation Under Bounded Context}
\label{sec:escalation}

% TODO: From Paper 2:
% - The delegation overhead definition
% - Theorem (Convergence to Escalation at Scale)
%   with the crossover at k*alpha >= C_n/2
% - The one-directionality argument
% - The circuit breaker pattern
% - Cell division illustration (condensed)

% PLACEHOLDER: Escalation section to be assembled from Paper 2.

% ==================================================================
% SECTION 5: NORMALIZATION
% ==================================================================
\section{Knowledge as Normalization}
\label{sec:normalization}

% TODO: From Paper 3:
% - Theorem (Normalization of Shared Propositions)
%   with the O(n*k) -> O(n+k') reduction
% - Corollary: Retention Criterion
% - Corollary: K/A Inseparability
% - The Codd analogy
% - Catastrophic forgetting as empirical signature

% PLACEHOLDER: Normalization section to be assembled from Paper 3.

% ==================================================================
% SECTION 6: PERSISTENCE AND INTELLIGENCE
% ==================================================================
\section{What Persists: Structural Invariants and Intelligence}
\label{sec:intelligence}

% TODO: From Paper 4:
% - Theorem (Evaluation-Driven State Revision / monotone descent)
% - Encoding Permanence: theta_i = 1 - epsilon/C_i
%   (the cleanest standalone result)
% - Theorem (A/M Stratification)
% - Theorem (Structural Invariants of Persistence)
%   gamma > 0 and eta_M > 0 as necessary conditions
% - Corollary: I(E) = gamma * eta_M as a natural measure
% - The free-context exhaustion argument
% - Necessary but not sufficient acknowledgment

% PLACEHOLDER: Intelligence section to be assembled from Paper 4.

% ==================================================================
% SECTION 7: EMPIRICAL CONFIRMATION
% ==================================================================
\section{Empirical Confirmation: Cancer as Architecture Failure}
\label{sec:empirical}

% TODO: Brief summary of Papers 5-7 results:
% - Paper 5: Channel count predicts OS with p < 10^{-14};
%   mutation count does not (the Eddington eclipse)
% - Paper 6: Channel structure maps to cell/tissue/organ/organism tiers
%   with bottom-up and top-down per tier
% - Paper 7: Graph position within the channel predicts treatment response
% - Cancer as defection from escalation back to swarm
% - The evolutionary data: three epochs, top-down channels scaled
%   with vertebrate complexity

% PLACEHOLDER: Empirical section to be assembled.

% ==================================================================
% SECTION 8: DISCUSSION
% ==================================================================
\section{Discussion}
\label{sec:discussion}

% TODO: The unification claim:
% - One constraint (bounded context) predicts four architectural
%   consequences
% - These consequences are observed independently in software,
%   biology, organizations, and AI
% - The contribution is not any individual result but the unification
% - Falsifiable predictions (condensed from Papers 1 and 4)
% - The encoding permanence threshold as the cleanest standalone result
% - Implications for AI architecture: CP with CQRS, not AP alone
% - Scope: information-theoretic invariants, not complete list
%   of everything needed to survive

% PLACEHOLDER: Discussion to be written fresh for the merged paper.

% ==================================================================
% SECTION 9: RELATED WORK
% ==================================================================
\section{Related Work}
\label{sec:related}

% TODO: One unified related work section:
% - CAP theorem (Brewer, Gilbert/Lynch)
% - Bounded rationality (Simon)
% - Knowledge representation (Hogan, Pearl, Baader, Berners-Lee, Novak)
% - Information bottleneck (Tishby)
% - Codd's relational model
% - Ashby's requisite variety
% - Free Energy Principle (Friston)
% - Intelligence measures (Hutter, Legg, Chollet, Schmidhuber)
% - Williamson on transaction costs
% - Catastrophic forgetting (McCloskey, Kirkpatrick)

% PLACEHOLDER: Related work to be consolidated from all 5 papers.

% ==================================================================
% CONCLUSION
% ==================================================================
\section{Conclusion}
\label{sec:conclusion}

% TODO: The single-constraint argument:
% - Bounded context forces graph structure, escalation,
%   normalization, and selects for agency + learning efficiency
% - These are not design preferences but cost-optimal responses
% - The architecture is the same across domains because the
%   constraint is the same
% - Empirical confirmation from cancer genomics
% - The encoding permanence threshold as a bridge to evolutionary
%   biology (learning rate hierarchy derived from cost-of-failure)

% PLACEHOLDER: Conclusion to be written fresh.

\bibliographystyle{plain}
\bibliography{../shared/references}

\end{document}
